\subsection{Illustrative U.S.--China--rest-of-world mapping}
\label{app:usmap}

This mapping places a unilateral US tariff on China in the three-country model, with $A$ the United States, $B$ China and $C$ the rest of the world. It maps trade shares and country sizes; it is not a calibration of exchange-rate magnitudes.

All inputs are for 2024, in billions of current dollars. US goods imports are $3{,}266.4$ in total and $440.3$ from China, and US goods exports $2{,}061.7$ in total and $143.3$ to China (US Census Bureau). China's goods imports are $2{,}585.1$ in total and $163.6$ from the United States, and its exports $3{,}577.2$ in total and $524.7$ to the United States (China's General Administration of Customs). This gives $b_{AB}=0.135$ and $b_{BA}=0.063$. The rest of the world's imports from the United States are US exports less US exports to China, and its imports from China are China's exports less China's exports to the United States, which gives $b_{CA}=0.386$ and $b_{CB}=0.614$. Country sizes are GDP at purchasing-power parity (IMF World Economic Outlook, October 2026), with the rest of the world equal to world GDP less that of the United States and China; relative to the United States they are $(1,1.30,4.45)$. The origin weights $\beta_{ij}$ are chosen so that the balanced free-trade baseline, solved at $\bar\rho=3\rho_D=3.96$ and then held fixed as $\rho$ varies, reproduces these six shares exactly. The remaining parameters are $\alpha_T=0.4$ and $\eta=1.5$, with either a common $\alpha_D=0.8$ or, for the United States, $\alpha_D=0.673$, the domestic share of US goods absorption in 2024 (BEA goods output plus imports less exports).

The bilateral bystander threshold, above which the dollar depreciates against the rest of the world, is $4.01$ with the common home bias and $3.87$ with the measured US home share. The aggregate terms-of-trade condition \cref{eq:totcondition} fails only for $\rho$ above $4.66$ and $4.50$, respectively. MML holds at both baselines for all $\rho\in[1,20]$. The mapping imposes balanced trade on every country, whereas the observed US goods deficit in 2024 was about $1{,}205$; it ignores the gaps between mirror trade statistics, uses model-implied rather than measured home shares for China and the rest of the world, covers goods only and treats the rest of the world as a single country.

\subsection{Data notes for the 2025 episode}
\label{app:episode-data}

\paragraph{Exchange-rate data.} Bilateral rates are Federal Reserve H.10 noon buying rates, expressed in dollars per unit of foreign currency (the CNY, CAD and MXN quotes are inverted). The broad index is the H.10 nominal broad dollar index (goods and services), which rises when the dollar appreciates; it is inverted. The index excluding China removes the renminbi from the daily log change of the broad index using China's 2025 H.10 weight (10.9 percent) and divides by one minus that weight, so the remaining weights are renormalized to sum to one. No daily dong series is available from the H.10 release or from the BIS; the VND row of \cref{tab:phases} uses BIS monthly averages in the monthly columns and, in the daily columns (marked $^{a}$), BIS end-of-month rates for the same months (end-January to end-March, end-March to end-April, end-July to end-December). Sources: FRED (DEXUSEU, DEXCHUS, DEXCAUS, DEXMXUS, DTWEXBGS), Federal Reserve H.10 trade weights and BIS (VND). ``Monthly'' compares monthly averages, January with March (Phase~1), March with April (Phase~2) and July with December (Phase~3), windows that do not coincide with the phase dates; ``daily'' runs from January~31 to April~2, April~2 to May~12 and August~6 to December~31. Panel~(a) of \cref{fig:episode} plots monthly averages, at mid-month, relative to each series' January~2025 average; the August dong average is missing in the source.

\paragraph{Tariff data.} Panel~(b) of \cref{fig:episode} plots statutory headline tariffs above the pre-2025 baseline by effective date; Section~232 product tariffs, exclusions and de minimis changes are omitted, and rates are not trade-weighted. Effective dates (2025), in percent: US on China 10 (February~4), 20 (March~4), 30 (April~5), 104 (April~9), 145 (April~10), 30 (May~14), 20 (November~10); China on US, across the board, 84 (April~10), 125 (April~12), 10 (May~14), with the product-specific measures of February~10 and March~10 not shown; US on EU 10 (April~5), 20 (April~9 only), 10 (April~10), 15 inclusive of MFN duties (August~7); US baseline on unlisted partners 10 (April~5); US on Canada 25 (March~4), non-USMCA goods only from March~7, 35 (August~1), with the rate on Mexico at 25 from March~4 to year-end (not shown); Canada on US 25 on covered goods (March~4), with coverage widened on March~13 and April~9 and narrowed to steel, aluminium and autos on September~1. US measures are from Federal Register executive orders and Canadian measures from the Department of Finance Canada; Chinese measures are as reported in USDA Foreign Agricultural Service summaries of State Council Tariff Commission announcements.

\paragraph{Phase definitions.} Phase~1 (February~1--April~2) covers the two US tariff increases on China and China's product-specific tariffs of 10--15 percent, on coal, liquefied natural gas, crude oil and some machinery and vehicles from February~10 and on agricultural products from March~10. In Phase~2 (April~3--May~12) the 10 percent baseline took effect on April~5, with Canada and Mexico excluded, and higher country-specific rates took effect on April~9 and were suspended for every partner except China on April~10. Phase~3 (August~7--December~31) begins with the August framework rates. Between the May~12 and August~6 daily observations, which fall outside the phases, the dollar fell by a further $4.9$ percent against the euro and $2.5$ percent against the broad index; in December it was $13.1$ percent weaker against the euro than in January.

\paragraph{Classification caveats.} The US tariffs on Canada (10 percent on energy) and Mexico took effect on March~4 and from March~7 no longer applied to goods qualifying under the USMCA. Canada retaliated on March~4 with a 25 percent surtax on C\$30 billion of US goods \citep{canadafinance2025} and extended it to a further C\$29.8 billion on March~13; Mexico did not retaliate. Canada and Mexico together carry 27.6 percent of the broad index. In the week of March~3--7 the broad index fell by $1.8$ percent, more than its net change over Phase~1, and the dollar fell by $4.4$ percent against the euro. The same week also contained the March~4 announcement by Germany's prospective governing parties of a large defence and infrastructure package and a loosening of the debt brake \citep{ripper2025}. The dollar strengthened against the dong by $3.2$ percent between the January and December averages, all of it by July. Vietnam, however, moved during the year from a bystander-like position to a direct target: a 46 percent rate was announced on April~2 and was in force for one day before the suspension of April~10, and a 20 percent rate applied from August~7. From August~7 to November~10, Vietnam's 20 percent against China's 30 percent was relative treatment in Vietnam's favour; after November~10, when both faced 20 percent, it disappeared. Within Phase~3 the dollar was roughly flat against the dong ($0.0$ percent on monthly averages, $+0.5$ percent between end-of-month rates).
